\documentclass[11pt]{article}
\usepackage[margin=1in]{geometry}
\usepackage{amsmath,amssymb,mathtools}
\usepackage{bm}
\usepackage{booktabs}
\usepackage{hyperref}

\title{Two-Stage Offset-Free Target Selector Report}
\author{}
\date{\today}

\begin{document}
\maketitle

\section{Purpose}
This note documents the implementation in
\texttt{Lyapunov/target\_selector.py}, specifically the function
\texttt{compute\_ss\_target\_refined\_rawlings(...)}.
The selector computes a steady-state target for offset-free tracking MPC
using a two-stage optimization:
\begin{enumerate}
    \item an exact offset-free target solve, and
    \item a closest-feasible fallback solve if the exact target is not admissible.
\end{enumerate}

\section{Model and Assumptions}
The code assumes the augmented linear model has the form
\begin{align}
    x_{k+1}^{\mathrm{aug}} &= A_{\mathrm{aug}} x_k^{\mathrm{aug}} + B_{\mathrm{aug}} u_k, \\
    y_k &= C_{\mathrm{aug}} x_k^{\mathrm{aug}},
\end{align}
with augmented state ordered as
\begin{align}
    x^{\mathrm{aug}} = \begin{bmatrix} x \\ d \end{bmatrix},
\end{align}
where:
\begin{align}
    x &\in \mathbb{R}^{n_x} && \text{physical state}, \\
    d &\in \mathbb{R}^{n_y} && \text{constant disturbance estimate}, \\
    u &\in \mathbb{R}^{n_u} && \text{control input}, \\
    y &\in \mathbb{R}^{n_y} && \text{controlled output.}
\end{align}

The matrices are partitioned in the implementation as
\begin{align}
    A &= A_{\mathrm{aug}}[1{:}n_x,1{:}n_x], \\
    B_d &= A_{\mathrm{aug}}[1{:}n_x,n_x{+}1{:}n_x{+}n_y], \\
    B &= B_{\mathrm{aug}}[1{:}n_x,:], \\
    C &= C_{\mathrm{aug}}[:,1{:}n_x], \\
    C_d &= C_{\mathrm{aug}}[:,n_x{+}1{:}n_x{+}n_y].
\end{align}
In mathematical notation this corresponds to
\begin{align}
    x_{k+1} &= A x_k + B u_k + B_d d_k, \\
    d_{k+1} &= d_k, \\
    y_k &= C x_k + C_d d_k.
\end{align}

The disturbance target is frozen at the current estimate:
\begin{align}
    d_s = \hat d.
\end{align}

\section{Steady-State Equations}
At steady state, the target variables $(x_s,u_s)$ satisfy
\begin{align}
    (I - A)x_s - B u_s - B_d \hat d = 0.
\end{align}
The associated steady-state output is
\begin{align}
    y_s = C x_s + C_d \hat d.
\end{align}

If all outputs are targeted, the controlled-output expression is
\begin{align}
    y_c = y_s.
\end{align}
If only selected output combinations are targeted, the optional matrix
$H \in \mathbb{R}^{n_c \times n_y}$ is used:
\begin{align}
    y_c = H y_s.
\end{align}
The requested steady-state target is denoted $y_{\mathrm{sp}}$.

\section{Bounds and Tightening}
The code uses tightened input bounds
\begin{align}
    u_{\ell} &= u_{\min} + u_{\mathrm{tight}}, \\
    u_{h} &= u_{\max} - u_{\mathrm{tight}}.
\end{align}
If output bounds are provided, the tightened output limits are
\begin{align}
    y_{\ell} &= y_{\min} + y_{\mathrm{tight}}, \\
    y_{h} &= y_{\max} - y_{\mathrm{tight}}.
\end{align}

\section{Stage 1: Exact Offset-Free Target Solve}
The first stage solves for an exact admissible equilibrium target. The
optimization problem is
\begin{align}
    \min_{x_s,u_s} \quad & (u_s-u_{\mathrm{nom}})^\top R_u (u_s-u_{\mathrm{nom}})
    + x_s^\top Q_x x_s
\end{align}
subject to
\begin{align}
    (I-A)x_s - B u_s - B_d \hat d &= 0, \\
    y_c &= y_{\mathrm{sp}}, \\
    u_{\ell} \le u_s &\le u_h.
\end{align}

If hard output bounds are enabled, the additional constraints are
\begin{align}
    y_{\ell} \le y_s \le y_h.
\end{align}

If soft output bounds are enabled, only the output inequalities are softened:
\begin{align}
    y_s + s_{\ell} &\ge y_{\ell}, \\
    y_s - s_{h} &\le y_h, \\
    s_{\ell} &\ge 0, \\
    s_{h} &\ge 0,
\end{align}
with extra penalties
\begin{align}
    s_{\ell}^\top W_{\ell} s_{\ell}
    \quad \text{and} \quad
    s_{h}^\top W_{h} s_{h}.
\end{align}
The exact equality constraint $y_c = y_{\mathrm{sp}}$ is never softened in Stage 1.

\section{Stage 2: Closest-Feasible Fallback}
If Stage 1 fails because the target is infeasible or numerically unacceptable,
the code solves a fallback problem:
\begin{align}
    \min_{x_s,u_s} \quad
    & (y_c-y_{\mathrm{sp}})^\top T_y (y_c-y_{\mathrm{sp}})
    + (u_s-u_{\mathrm{nom}})^\top R_u (u_s-u_{\mathrm{nom}})
    + x_s^\top Q_x x_s.
\end{align}
If previous target information is provided, the fallback stage also includes
\begin{align}
    (x_s-x_s^{\mathrm{prev}})^\top Q_{dx}(x_s-x_s^{\mathrm{prev}})
    + (u_s-u_s^{\mathrm{prev}})^\top R_{du}(u_s-u_s^{\mathrm{prev}}).
\end{align}

The fallback constraints are
\begin{align}
    (I-A)x_s - B u_s - B_d \hat d &= 0, \\
    u_{\ell} \le u_s &\le u_h,
\end{align}
plus the same hard or soft output-bound constraints, if provided. There is no
hard equality $y_c = y_{\mathrm{sp}}$ in Stage 2.

\section{Acceptance Checks}
The implementation does not accept a solution only because the solver reports
\texttt{optimal} or \texttt{optimal\_inaccurate}. It performs explicit
post-solve checks.

The dynamic residual is
\begin{align}
    r_{\mathrm{dyn}} = (I-A)x_s - B u_s - B_d \hat d.
\end{align}
Its infinity norm is
\begin{align}
    \|r_{\mathrm{dyn}}\|_{\infty}.
\end{align}

For Stage 1, the target equality residual is
\begin{align}
    r_{\mathrm{eq}} = y_c - y_{\mathrm{sp}},
\end{align}
with infinity norm $\|r_{\mathrm{eq}}\|_{\infty}$.

The code uses the following acceptance tolerances:
\begin{align}
    \|r_{\mathrm{dyn}}\|_{\infty} &\le 10^{-6}
    && \text{if status is \texttt{optimal}}, \\
    \|r_{\mathrm{dyn}}\|_{\infty} &\le 10^{-5}
    && \text{if status is \texttt{optimal\_inaccurate}}.
\end{align}
For Stage 1 the same tolerance rule is applied to $\|r_{\mathrm{eq}}\|_{\infty}$.

If output bounds are hard, the code also checks the bound violation magnitude
numerically and rejects the solution if the violation exceeds the corresponding
tolerance.

\section{Returned Quantities}
On success, the selector returns
\begin{align}
    x_s, \qquad u_s, \qquad d_s = \hat d.
\end{align}
It also constructs the augmented steady-state target
\begin{align}
    x_s^{\mathrm{aug}} = \begin{bmatrix} x_s \\ \hat d \end{bmatrix}.
\end{align}

The debug dictionary includes:
\begin{align}
    \texttt{solve\_stage} &\in \{\texttt{"exact"},\texttt{"fallback"}\}, \\
    \texttt{target\_error} &= y_c - y_{\mathrm{sp}}, \\
    \texttt{target\_slack} &= \texttt{target\_error}, \\
    \texttt{target\_eq\_residual\_inf} &= \|r_{\mathrm{eq}}\|_{\infty}, \\
    \texttt{dyn\_residual\_inf} &= \|r_{\mathrm{dyn}}\|_{\infty}.
\end{align}

The alias
\begin{align}
    \texttt{target\_slack} = \texttt{target\_error}
\end{align}
is retained for backward compatibility with downstream diagnostics and exporters.

\section{Interpretation}
The two-stage structure matters because it separates two different objectives:
\begin{enumerate}
    \item finding a true equilibrium target that exactly matches the requested setpoint when feasible, and
    \item finding the nearest admissible equilibrium when exact tracking is impossible.
\end{enumerate}
This is the right structure for Lyapunov tracking MPC, because the target center
should be the true equilibrium whenever an exact equilibrium exists.

\section{Python-to-Math Mapping}
\begin{center}
\begin{tabular}{ll}
\toprule
Python name & Mathematical meaning \\
\midrule
\texttt{A\_aug, B\_aug, C\_aug} & augmented linear model \\
\texttt{xhat\_aug} & current augmented estimate $[\hat x;\hat d]$ \\
\texttt{d\_hat} & frozen disturbance estimate $\hat d$ \\
\texttt{y\_sp} & requested steady-state target $y_{\mathrm{sp}}$ \\
\texttt{u\_nom} & nominal input center $u_{\mathrm{nom}}$ \\
\texttt{Ty\_diag} & diagonal of $T_y$ \\
\texttt{Ru\_diag} & diagonal of $R_u$ \\
\texttt{Qx\_diag} & diagonal of $Q_x$ \\
\texttt{Qdx\_diag} & diagonal of $Q_{dx}$ \\
\texttt{Rdu\_diag} & diagonal of $R_{du}$ \\
\texttt{Wy\_low\_diag, Wy\_high\_diag} & diagonals of $W_{\ell}, W_h$ \\
\texttt{H} & controlled-output selector matrix \\
\bottomrule
\end{tabular}
\end{center}

\section{Recommended Use in a Written Report}
If this material is going into a thesis note, project report, or paper draft, the
best structure is:
\begin{enumerate}
    \item describe the augmented offset-free model,
    \item state the steady-state equations,
    \item present Stage 1 as the exact target selector,
    \item present Stage 2 as the artificial-reference fallback,
    \item explain the residual-based acceptance rule,
    \item then connect the selected equilibrium target to the Lyapunov MPC layer.
\end{enumerate}

\end{document}
